% Part: sets-functions-relations
% Chapter: functions
% Section: inverses

\documentclass[../../../include/open-logic-section]{subfiles}

\begin{document}

\olfileid{sfr}{fun}{inv}
\olsection{Functies omkeren}

\begin{explain}
Je kunt een functie zien als iets dat een invoer naar een uitvoer
stuurt. Dan ligt één vraag voor de hand: kan dat ook ``omgekeerd''?
Bijvoorbeeld: de opvolgerfunctie $f(x) = x + 1$ kun je omkeren. De
functie $g(y) = y - 1$ maakt dan ``ongedaan'' wat $f$ doet.

Maar je moet wel opletten. De definitie van~$g$ geeft een functie
$\Int \to \Int$, maar geen \emph{functie} $\Nat \to \Nat$, want
$g(0) \notin \Nat$. Zelfs bij zo'n eenvoudig voorbeeld is dus niet
meteen duidelijk of je een functie kunt omkeren. Dat kan afhangen van
het domein en codomein: van de verzameling waaruit de invoer komt en de
verzameling waarin de uitvoer moet liggen.

Daarvoor gebruiken we het begrip inverse van een functie.
\end{explain}

\begin{defn}
Een functie $g \colon B \to A$ is een \emph{inverse} van een functie
$f \colon A \to B$ als $f(g(y)) = y$ en $g(f(x)) = x$ voor alle $x
\in A$ en $y \in B$.
\end{defn}

Als $f$ een inverse~$g$ heeft, schrijven we vaak $f^{-1}$ in plaats
van~$g$.

\begin{explain}
Nu zoeken we uit wanneer een functie een inverse heeft. Voor
$f\colon A \to B$ lijkt $g\colon B \to A$ een goede keuze. We zouden
deze functie als volgt willen ``definiëren'':
\[
g(y) = \text{``de'' $x$ waarvoor $f(x) = y$.}
\]
De aanhalingstekens rond ``definiëren'' (en rond ``de'') verraden al
dat dit nog geen echte definitie is. Dit werkt namelijk niet altijd.
Om hiermee een functie te bepalen, moet er precies één~$x$ zijn
waarvoor $f(x) = y$. De uitvoer van~$g$ moet maar op één manier zijn
vastgelegd. Ook moet er voor elke $y \in B$ een uitvoer zijn. Als er
$x_1$ en $x_2 \in A$ zijn met $x_1 \neq x_2$ maar $f(x_1) = f(x_2)$,
dan heeft $g(y)$ geen unieke waarde voor $y = f(x_1) = f(x_2)$. En als
er helemaal geen~$x$ is waarvoor $f(x) = y$, dan is $g(y)$ helemaal
niet vastgelegd. Kortom, $g$ is alleen zo te definiëren als $f$ zowel
!!{injective} is (verschillende invoeren krijgen verschillende
uitvoeren) als !!{surjective} (elke mogelijke uitvoer wordt bereikt).

We nemen dit stap voor stap en splitsen de vraag in twee delen. Neem
een functie~$f\colon A \to B$. Wanneer bestaat er een functie
$g\colon B \to A$ waarvoor $g(f(x)) = x$? Zo'n $g$ ``maakt ongedaan''
wat $f$ doet en heet een \emph{linksinverse} van~$f$. En wanneer bestaat
er een functie $h\colon B \to A$ waarvoor $f(h(y)) = y$? Zo'n $h$
heet een \emph{rechtsinverse} van~$f$: $f$ ``maakt ongedaan'' wat
$h$~doet.
\end{explain}

\begin{prop}
Als het domein niet leeg is en $f\colon A \to B$ !!{injective} is, dan
bestaat er een
\emph{linksinverse}~$g\colon B \to A$ van~$f$ waarvoor $g(f(x)) = x$
voor alle $x \in A$.
\end{prop}

\begin{proof}
Stel dat het domein niet leeg is en dat $f\colon A \to B$
!!{injective} is. Neem een $y \in B$. Als
$y \in \ran{f}$, dan is er een $x \in A$ waarvoor $f(x) = y$. Omdat
$f$ !!{injective} is, is er maar één zo'n~$x \in A$. We kunnen dus
$g(y) = x$ nemen. Anders gezegd: $g(y)$ is ``de'' $x \in A$ waarvoor
$f(x) = y$. Als $y \notin \ran{f}$, mogen we die naar elk
willekeurig~$a \in A$ sturen. We kiezen daarom een $a \in A$ en leggen
$g\colon B \to A$ zo vast:
\[
g(y) = \begin{cases}
    x & \text{als $f(x) = y$}\\
    a & \text{als $y \notin \ran{f}$.}
\end{cases}
\]
De functie is nu voor elke $y \in B$ vastgelegd: als zo'n $y \in \ran{f}$
ligt, is er precies één $x \in A$ waarvoor $f(x) = y$. Uit de
definitie volgt dat als $y = f(x)$, dan $g(y) = x$, en dus
$g(f(x)) = x$.
\end{proof}

\begin{prob}
Laat zien dat als $f\colon A \to B$ een linksinverse~$g$ heeft,
$f$~!!{injective} is.
\end{prob}

\begin{prop}
    Als $f\colon A \to B$ !!{surjective} is, dan bestaat er een
    \emph{rechtsinverse}~$h\colon B \to A$ van~$f$ waarvoor $f(h(y)) =
    y$ voor alle~$y \in B$.
\end{prop}

\begin{proof}
Stel dat $f\colon A \to B$ !!{surjective} is. Neem een $y \in B$.
Omdat $f$~!!{surjective} is, is er een $x_y \in A$ met $f(x_y) = y$.
We kunnen dus $h(y) = x_y$ nemen. Dat betekent: voor elke $y \in B$
kiezen we een $x \in A$ waarvoor $f(x) = y$. Omdat $f$~!!{surjective}
is, valt er altijd minstens één te kiezen.\footnote{Omdat $f$
!!{surjective} is, is voor elke~$y \in B$ de verzameling
$\Setabs{x}{f(x) = y}$ niet leeg. Voor onze definitie van~$h$ moeten
we uit elk van deze verzamelingen één $x$ kiezen. Het spreekt niet
vanzelf dat dit altijd kan. We nemen die mogelijkheid gewoon als een
axioma aan. Anders gezegd: deze stelling gebruikt het zogeheten
keuzeaxioma. Dit is een punt waarop we
\oliflabeldef{sth:choice::chap}{terugkomen in \olref[sth][choice][]{chap}}{hier niet verder ingaan}.
Maar in veel concrete gevallen is het keuzeaxioma niet nodig,
bijvoorbeeld als $A = \Nat$ of eindig is, of als $f$ !!{bijective} is.
(Als $f$ !!{bijective} is, heeft voor elke $y \in B$ de verzameling
$\Setabs{x}{f(x) = y}$ precies één !!{element}.
Er valt dan dus niets te kiezen.)}
Volgens de definitie geldt: als $x = h(y)$, dan $f(x) = y$. Dus voor
elke $y \in B$ geldt $f(h(y)) = y$.
\end{proof}

\begin{prob}
Laat zien dat als $f\colon A \to B$ een rechtsinverse~$h$ heeft,
$f$~!!{surjective} is.
\end{prob}

\begin{explain}
  Als we de ideeën uit het vorige bewijs samenvoegen, zien we dat elke
  !!{bijection} een inverse heeft. Er is dus één functie die tegelijk
  linksinverse en rechtsinverse van~$f$ is.
\end{explain}

\begin{prop}\ollabel{prop:bijection-inverse}
Als $f\colon A \to B$ !!{bijective} is, dan bestaat er een
functie~$f^{-1}\colon B \to A$ waarvoor voor alle $x \in A$ geldt:
$f^{-1}(f(x)) = x$, en voor alle $y \in B$: $f(f^{-1}(y)) = y$.
\end{prop}

\begin{proof}
Oefening.
\end{proof}

\begin{prob}
Bewijs \olref[sfr][fun][inv]{prop:bijection-inverse}. Je moet~$f^{-1}$
definiëren, laten zien dat dit een functie is en laten zien
dat het een inverse van~$f$ is. Dus: $f^{-1}(f(x)) = x$ en
$f(f^{-1}(y)) = y$ voor alle $x \in A$ en $y \in B$.
\end{prob}

\begin{explain}
Je kunt ook in iets meer gevallen aan een inverse komen. In
\olref[kin]{sec} zagen we dat elke functie $f$ een !!a{surjection} $f'
\colon A \to \ran{f}$ oplevert door $f'(x) = f(x)$ te nemen voor alle
$x \in A$. Als $f$~!!{injective} is, dan is $f'$~!!{bijective}.
Volgens \olref{prop:bijection-inverse} heeft deze functie dus een
unieke inverse. We zijn dan een klein beetje slordig met de notatie en
noemen de inverse van $f'$ soms gewoon ``de inverse van~$f$.''
\end{explain}

\begin{prop}\ollabel{prop:left-right}%
  Laat zien dat als $f\colon A \to B$ een linksinverse~$g$ en een
  rechtsinverse~$h$ heeft, dan $h = g$.
\end{prop}

\begin{proof}
  Oefening.
\end{proof}

\begin{prob}
  Bewijs \olref[sfr][fun][inv]{prop:left-right}.
\end{prob}

\begin{prop}\ollabel{prop:inverse-unique}
Een functie~$f$ kan nooit meer dan één inverse hebben.
\end{prop}

\begin{proof}
  Stel dat $g$ en $h$ allebei inversen van~$f$ zijn. Dan is $g$ in elk
  geval een linksinverse van~$f$ en is $h$ een rechtsinverse. Uit
  \olref{prop:left-right} volgt $g = h$.
\end{proof}
\end{document}
